% Part: first-order-logic
% Chapter: sequent-calculus
% Section: structural-rules

\documentclass[../../../include/open-logic-section]{subfiles}

\begin{document}

\iftag{FOL}
      {\olfileid{fol}{seq}{srl}}
      {\olfileid{pl}{seq}{srl}}

\olsection{গঠনগত বিধি}

সিকোয়েন্টের বাঁ ও ডান পাশে !!{sentence}s পুনর্বিন্যাস করার জন্যও কয়েকটি
বিধি দরকার। কোনো যৌক্তিক বিধি যে পূর্বধারণার !!{sentence}s-এর উপর কাজ
করে, সেগুলিকে একেবারে বাঁ বা ডান প্রান্তে থাকতে হয়। তাই !!{sentence}s-কে
সঠিক স্থানে সরানোর জন্য একটি “অদলবদল” বিধি দরকার। কখনো দুটি অভিন্ন
!!{sentence}s-কে একটিতে মেলানো এবং যেকোনো পাশে একটি !!a{sentence} যোগ
করতে পারাও গুরুত্বপূর্ণ।

\subsection{দুর্বলীকরণ}

\begin{defish}
\Axiom$ \Gamma \fCenter \Delta $
\RightLabel{\LeftR{\Weakening}}
\UnaryInf$ !A, \Gamma \fCenter \Delta$
\DisplayProof
\hfill
\Axiom$ \Gamma \fCenter \Delta$
\RightLabel{\RightR{\Weakening}}
\UnaryInf$ \Gamma \fCenter \Delta, !A$
\DisplayProof
\end{defish}

\subsection{সংকোচন}

\begin{defish}
\Axiom$ !A, !A, \Gamma \fCenter \Delta $
\RightLabel{\LeftR{\Contraction}}
\UnaryInf$ !A, \Gamma \fCenter \Delta$
\DisplayProof
\hfill
\Axiom$ \Gamma \fCenter \Delta, !A, !A$
\RightLabel{\RightR{\Contraction}}
\UnaryInf$ \Gamma \fCenter \Delta, !A$
\DisplayProof
\end{defish}

\subsection{অদলবদল}

\begin{defish}
\Axiom$ \Gamma, !A, !B, \Pi \fCenter \Delta $
\RightLabel{\LeftR{\Exchange}}
\UnaryInf$ \Gamma, !B, !A, \Pi \fCenter \Delta$
\DisplayProof
\hfill
\Axiom$ \Gamma \fCenter \Delta, !A, !B, \Lambda$
\RightLabel{\RightR{\Exchange}}
\UnaryInf$ \Gamma \fCenter \Delta, !B, !A, \Lambda$
\DisplayProof
\end{defish}

দুর্বলীকরণ, সংকোচন ও অদলবদল অনুমানের কোনো ধারাকে প্রায়ই দ্বিগুণ
অনুমান-রেখা দিয়ে দেখানো হবে।

নিচের বিধিটির নাম “কর্তন”; কঠোর অর্থে এটি আবশ্যক নয়, তবে !!{derivation}s
পুনর্ব্যবহার ও একত্র করা অনেক সহজ করে।
\begin{defish}
\[
\Axiom$ \Gamma \fCenter \Delta, !A$
\Axiom$ !A, \Pi \fCenter \Lambda $
\RightLabel{\Cut}
\BinaryInf$ \Gamma, \Pi \fCenter \Delta, \Lambda$
\DisplayProof
\]
\end{defish}

\end{document}
